\subsection{Five states over a three-by-two minimal band}

\begin{lemma}[Invariant small column sets in the remaining three-by-five hull]
\label{C:3x5-small-invariant-columns}
Let $H$ be a local-character-allowed actor submonoid of orbit $164$
in Theorem~\ref{B:three-by-five-targets}, containing its full
$3\times5$ band and identity.  Suppose that a nonempty subset of
at most two columns is invariant under every actor outside the
band.  Then every product of rational degree-two classes of the
ambient orbit-$164$ hull restricts to zero on $BH$.
More precisely, $H$ lies in one of five
explicit actor monoids with zero fourth rational homology.
\end{lemma}
\begin{proof}
The allowed maximal actor monoid has thirty elements, consisting
of the fifteen band elements, its identity, and fourteen optional
actors.  Test all $2^{14}$ optional subsets for multiplication
closure and for invariance of each of the five singleton and ten
two-element column subsets.  The complete counts are
\[
 \begin{array}{l|r}
 \text{closed actor monoids}&2186\\
 \text{closed monoids admitting such an invariant subset}&408\\
 \text{maximal monoids in the latter family}&5.
 \end{array}
\]
In the indexing of the complete safe-submonoid list, the five
maxima and their independently certified bar ranks are
\[
 \begin{array}{c|r|r|r|r}
 \text{index}&\text{order}&\operatorname{rank}d_2&
 \operatorname{rank}d_3&\dim C_2\\\hline
 1031&22&39&249&288\\
 1732&24&57&455&512\\
 1794&24&57&455&512\\
 1804&24&57&455&512\\
 1860&24&57&455&512.
 \end{array}
\]
The ranks are over $\mathbb F_5$ in the normalized two-sided
complex of Lemma~\ref{lem:B-band-Tor}; its degree two computes
fourth monoid homology.  In every case the two ranks sum to
the middle dimension, proving $H_4(-;\mathbb F_5)=0$, hence
$H^4(-;\mathbb Q)=0$.  Inclusions into these five monoids
therefore kill every pulled-back degree-two product.

The direct subset enumeration and every maximality comparison are
in \path{code/a_five_state_band_obstruction.py}.  It compares its
complete list with the independently certified orbit-$164$ list,
but does not use that list to generate subsets.  The output
\path{code/outputs/a_five_state_band_obstruction.json} includes
all closed sets, all $408$ retained images, and independent-column
witnesses for the five pairs of differentials.  A separate dense
table-based reconstruction, using the verifier of the earlier
three-by-three theorem with shape parameters $(3,5)$, checks
those rank witnesses and all associativity equations.
\end{proof}

\begin{theorem}[A five-state three-by-two obstruction]
\label{C:five-state-three-by-two-cones}
Let $M$ be a finite monoid with rectangular-band minimal ideal
of shape $3\times2$, and let $S$ be a right $M$-set with five
elements.  Then
\[
 B\mathcal C(M,S)\not\simeq\mathbb CP^2.
\]
There is no bound on the number of elements of $M$ in this statement.
\end{theorem}
\begin{proof}
Write $K=I\times J$, with $|I|=3$ and $|J|=2$, and suppose
that $X=B\mathcal C(M,S)$ has the indicated finite homotopy type.
The act is connected, since the act-cone homology sequence and
$H_1(BM;\mathbb Q)=0$ identify $H_1(X;\mathbb Q)$ with the
reduced degree-zero homology of the action category.
If a state is fixed by all of $K$, the first paragraph of
Lemma~\ref{C:four-state-band-dichotomy} applies without any
state-count restriction: the band grid maps onto $H_2(BM)$
through the edge space, and every degree-two product on $X$
is zero.  We may exclude this case.

First suppose that some element of $K$ acts nonconstantly.
The set $SK$ is a disjoint union of right $K$-orbits, each
invariant under $M$.  Each orbit has at least two elements and
at most two elements: its elements are obtained by varying the
column of a fixed band row.  If there were a single orbit,
idempotence would make every $K$-actor constant on all of $S$,
as in Lemma~\ref{C:four-state-band-dichotomy}.  Thus $SK$
consists of two two-point orbits, with the remaining state
outside $SK$.  Label the two orbits by their two column values.

Let $\bar S$ be the quotient in
Lemma~\ref{C:constant-action-quotient-H2}.  For either column,
its state in the first orbit is identified with its state in
the second orbit.  Therefore $|\bar S|\leq3$.
All elements of $K$ act constantly on $\bar S$, and that lemma
gives a surjection
\[
 H^2(B\mathcal C(M,\bar S);\mathbb Q)
       \longrightarrow H^2(X;\mathbb Q).
\]
Compress the quotient act cone by
Theorem~\ref{C:constant-minimal-act-compression}, obtaining
a monoid $N$ with a $3\times|\bar S|$ minimal band.
For at most two columns every balanced matrix has rank at most
one.  Lemma~\ref{B:rank-one-square} and polarization therefore
make the rational degree-two products zero, without requiring
finite homotopy type of $BN$.

For three columns, the proof of
Theorem~\ref{B:three-by-three-cups} applies to every class
pulled back to $X$.  Here is the finite-target point that is
needed for that assertion.  All cases other than its exceptional
skew-form orbit have universal cup primitives or zero fourth
cohomology witnesses, which pull back directly.  In the skew
case, an eventual local group of $N$ that maps onto a forbidden
cyclic actor group is nontrivial and therefore lies outside the
new band.  The outside elements and all products within such
a group are literally unchanged from $M$.  This supplies a
local-group map into the endomorphism monoid of
$\mathcal C(M,S)$.  The actual functor
\[
 \mathcal C(M,S)\longrightarrow\mathcal C(M,\bar S)
       \longrightarrow\operatorname{Kar}(N)
\]
pulls the integral hull cocycle back along this same group map.
Its forbidden cyclic character remains nonzero under inflation,
whereas Miller's theorem makes the map from that finite group's
classifying space to the finite CW model of $X$ null-homotopic.
This is the same contradiction as in the original skew-orbit
proof.  Thus its bad actors are excluded here as well, and the
remaining nine targets have zero fourth rational cohomology.
No finite-type assertion about $BN$ is used.
All pulled-back degree-two products consequently vanish.
Their degree-two surjection onto $X$ contradicts the square in
$H^*(\mathbb CP^2;\mathbb Q)$.

It remains that every element of $K$ acts constantly on $S$.
Compression now gives an actual homotopy equivalence
\[
 X\simeq BN,
\]
where $N$ has minimal band $I\times S$ of shape $3\times5$.
For $j\in J$, let $c_j\in S$ be the constant value of
any band element with column $j$.  The nonempty set
\[
 T=\{c_j:j\in J\},\qquad |T|\leq2,
\]
is invariant under every outside actor.  Indeed
$c_jm=c_{km}$ for any $k$ in column $j$, and $km\in K$.
Theorem~\ref{B:three-by-five-targets} reduces a possible
projective-plane source of this shape to its allowed orbit-$164$
action image.  That image fixes the entire band, and every
outside actor in the image comes from an outside element of
$N$; it therefore preserves the corresponding subset $T$.
Lemma~\ref{C:3x5-small-invariant-columns} makes the restricted
ambient degree-two square zero.  All rational degree-two classes
of $BN$ pull back from the ambient orbit-$164$ hull: the source,
its action image, and that maximal balanced-form hull have the
same rational balancing relation space.  Thus the
balanced-tensor presentation identifies their degree-two groups,
so this is again a contradiction.
\end{proof}

\begin{remark}[Scope and computational controls]
\label{C:five-state-band-cones-scope}
The preceding theorem assumes a minimal band, not a nontrivial
minimal maximal group.  It uses complete finite actor calculations
only in the already classified three-by-three and three-by-five
hulls and in the five new zero-homology witnesses above; it does
not enumerate monoids of unbounded order.
As a separate control, all $1082$ two-core-orbit five-state normal
forms over the $272$ nonsplit-row order-ten residual targets
were checked in \path{code/a_five_state_two_orbits.py}.
Of these, $1050$ fail exact rational homology through degree
three and the other $32$ have explicit mod-two square primitives.
Three complete independent action enumerations agree with the
normal-form generator.  These controls are not used to establish
the all-order theorem.
\end{remark}
